\section{The coordinate process and conditional revelation}\label{sec:prelim}

Write $G=\Sym(V)\cong S_n$. Permutations map labels to positions,
and $(gh)(x)=g(h(x))$.
For independent group-valued variables $g\sim\mu$, $h\sim\nu$,
their product has law $\mu*\nu$. Let
$R(x_1,\ldots,x_d)=(x_2,\ldots,x_d,x_1)$. One physical shuffle
is $RS$, where $S$ independently fixes or swaps every pair in
direction $e_1$. If $Y_t$ is the physical permutation, then
\begin{equation}\label{eq:frames:cyclic}
 \Pi_t=R^{-t}Y_t,\qquad
 \Pi_{t+1}=Q_t\Pi_t,\qquad i_t=1+(t\bmod d),
\end{equation}
where the $Q_t$ are independent fair matching switches in direction
$e_{i_t}$. At multiples of $d$, the moving frame is the identity.
Independent blocks of whole sweeps therefore have their physical
shuffle convolution laws. A deterministic change of starting deck
right-translates each law, preserving distance from uniform:
\begin{equation}\label{eq:worst-state}
 \sup_{g_0}\|\mathcal L(Y_t\mid Y_0=g_0)-U_{S_n}\|_{\TV}
 =\|q_d^{*t}-U_{S_n}\|_{\TV}.
\end{equation}
We use counting measure for vector norms, natural logarithms unless
marked otherwise, and unnormalized trace for Hilbert--Schmidt norms.
For probabilities $p,q$ on a finite set, relative entropy is
$D(p\Vert q)=\sum_xp(x)\log(p(x)/q(x))$, with $0\log0=0$
and value $+\infty$ if $p$ charges an atom on which $q$ vanishes.
The entropy of $p$ is $H(p)=-\sum_xp(x)\log p(x)$.

\begin{lemma}[Conditional path flow]\label{lem:marginal-averaging}
Fix observed labels and a different tagged label, independently of the
switches. Conditional on any feasible specification of the observed
paths through time $T$, the coins on edges not visited by those paths
remain independent and fair. The conditional tag law at time $s\le T$
is computed using the paths through $s$ only. It averages masses on
edges with two available positions and transports mass on an edge
with one available position to its uniquely available output. This
rule also carries uniform measure on the available set to uniform
measure on the next available set.
\end{lemma}

\begin{proof}
A feasible path specification prescribes exactly one coin value at
each visited time--edge pair. Necessity is immediate. Conversely,
chronological induction shows that those prescribed values force the
specified paths; no other coin is constrained. The event is thus a
cylinder in the independent coin array. Its restriction after time $s$
specifies only later time--edge coins: the observed positions at time
$s$ are already fixed by the path prefix. These later constraints are
independent of all earlier coins, so they do not alter the conditional
tag law at time $s$. On unvisited pairs, averaging
over a fresh fair coin gives the stated rule; visited pairs transport
the remaining endpoint. This proves the prefix-measurable update. The uniform vector has
equal masses on every available endpoint, so it obeys the same rule.
This is the path-flow principle of~\cite[Lemma~\ref*{h-lem:path-cylinder}]{optimalThorp}.
\end{proof}

\begin{lemma}[Sequential comparison]\label{lem:sequential-tv}\label{lem:sequential}
Let $(Z_1,\ldots,Z_k)$ be an injective random tuple in a finite set
$V$. Its distance from the uniform injective tuple is at most
\[
 \sum_{j=1}^k\E\left\|
 \mathcal L(Z_j\mid Z_1,\ldots,Z_{j-1})
 -\Unif(V\setminus\{Z_1,\ldots,Z_{j-1}\})
 \right\|_{\TV}.
\]
Each expectation uses the actual prefix law. It may be bounded by
conditioning further on earlier paths, provided the available endpoint
set is already determined by the prefix. More generally, let $\mathcal A$ be additional base information and let
$A\subseteq V$ be an $\mathcal A$-measurable set, with the tuple
supported on $\operatorname{inj}([k],A)$. The same statement holds
conditional on $\mathcal A$, replacing $V$ by $A$, and can then be
averaged over $\mathcal A$.
\end{lemma}

\begin{proof}
Compare laws with an actual prefix and a suffix filled by uniform
sampling without replacement. Two consecutive laws differ only in
the next conditional sampling kernel; adding their common suffix
cannot increase variation distance. Summing the differences gives
the displayed bound. Conditional convexity proves the full-path
and base-information variants. This is the comparison in
\cite[Lemma~\ref*{h-lem:sequential-tv}]{optimalThorp}.
The basic sequential inequality also appears in Morris's exclusion-process
analysis~\cite[Section~5.3, Lemma~12]{morris-exclusion2006} and in
Morris, Rogaway and Stegers~\cite[Lemma~2 and Appendix~A]{mrs}.
\end{proof}

For later use, the elementary support obstruction is
\begin{equation}\label{eq:support}
 \|q_d^{*t}-U_{S_n}\|_{\TV}\ge1-\frac{2^{tn/2}}{n!}
 \qquad(t\in\mathbb Z_{\ge0}).
\end{equation}
Indeed $t$ physical shuffles use $tn/2$ fair bits, so their support
has at most $2^{tn/2}$ elements. Uniform measure assigns it at most
that many divided by $n!$. Consequently the threshold-$1/4$ mixing
time is at least
$\lceil(2/n)\log_2(3n!/4)\rceil=2d-O(1)$. More explicitly,
$n!\ge(n/e)^n$ gives
\begin{equation}\label{mart:lower-exact}
 t_{\mathrm{mix}}(d)\ge
 2d-2\log_2 e+\frac2n\log_2(3/4).
\end{equation}
For $t\le d$ the distance is at least $1-(e/\sqrt n)^n$.
For fixed $d$ the chain
does converge: identity sweeps and any single cube-edge transposition
have positive probability; cube-edge transpositions generate $S_n$,
and padding with identity sweeps yields a common positive minorization.

\begin{lemma}[Concentration]\label{lem:concentration}
Suppose $F$ depends on independent variables, and changing variable $j$
changes $F$ by at most $c_j$. For $u\ge0$, either tail at distance $u$
from its mean has probability at most
$\exp(-2u^2/\sum_jc_j^2)$. For a martingale whose successive increments
have absolute values at most $c_j$, either tail is at most
$\exp(-u^2/(2\sum_jc_j^2))$. If all $c_j$ vanish the variable is constant.
\end{lemma}
\begin{proof}
For a random variable $Z$ in an interval of length $c$, the second
 derivative of $\log\E e^{\theta Z}$ is a tilted variance, at most
$c^2/4$. Integration twice gives
$\E e^{\theta(Z-\E Z)}\le e^{\theta^2c^2/8}$.
Expose the independent variables successively. The corresponding Doob
martingale difference has conditional range of length at most $c_j$:
two choices of the next variable change each possible conditional
completion by at most $c_j$. Multiplying the conditional exponential
bounds and optimizing in $\theta$ proves the first assertion. An
absolute increment bound $c_j$ instead gives conditional range of
length at most $2c_j$, proving the second. These are the exponential
arguments of Hoeffding and Azuma~\cite{hoeffding1963,azuma1967}.
\end{proof}

In particular, independent $[0,1]$ variables give either tail
$e^{-2u^2/r}$ for a sum of $r$ terms. For a uniform $m$-subset of
an even set, expose its points in random order. Coupling two possible
next draws by exchanging their identities in the remaining completion
shows that the Doob martingale for the final half-count has increments
bounded by one. Its mean is $m/2$, so the probability that its count
in a specified half is below $m/4$ is at most $e^{-m/32}$. These
concentration estimates will be applied before conditioning on full
observed paths.
